\documentclass[../../main.tex]{subfiles}
\begin{document}

\paragraph{【新課程】}
$a>1$とし，$u=a^{x+y}$，$v=a^{2x-y}$とおく．
$x$，$y$平面上の点$P(x,y)$が正方形$ABCD$の周に沿って一周するとき，
$u$，$v$平面上で対応する点$Q(u,v)$がえがく軌跡によって囲まれる図形の面積を求めよ．
ただし$A(1,1)$，$B(-1,1)$，$C(-1,-1)$，$D(1,-1)$とする．

\paragraph{【旧課程】}
$f(x)=x^3-6x^2+8$とし，$0 \leqq x \leqq r$における$|f(x)|$の最大値を$M(r)$とするとき，
積分$\displaystyle\int_0^5M(r)\,dr$を求めよ．

\end{document}
